% 05-turing-instability
\chapter{Two equations: diffusion-driven instability}\label{ch:turing}

We now come to the heart of these notes. In \cref{ch:one-equation} we saw that a
single reaction--diffusion equation cannot turn a stable constant state into an
unstable one: diffusion only adds damping. With two species the situation changes
completely. Our plan is the same as for the heat equation in \cref{ch:diffusion}:
linearize, expand in cosines, and observe that the linearized PDE splits into
infinitely many \emph{independent} linear ODE systems of size $2\times2$, one for
each cosine mode. The question ``can diffusion destabilize?'' then becomes a
question about $2\times 2$ matrices, which we answer completely.

\section{The setting}\label{sec:turing-setting}

We consider two chemical species with concentrations $u(x,t)$ and $v(x,t)$ on the
interval $(0,L)$:
\begin{equation}\label{eq:rd-system}
\begin{aligned}
  u_t &= d_u\, u_{xx} + f(u,v), & x\in(0,L),\ t>0,\\
  v_t &= d_v\, v_{xx} + g(u,v), & x\in(0,L),\ t>0,\\
  u_x &= v_x = 0, & x\in\{0,L\},\ t>0,\\
  u(x,0) &= u_0(x),\quad v(x,0)=v_0(x), & x\in(0,L).
\end{aligned}
\end{equation}
Here $d_u,d_v>0$ are the diffusion coefficients, collected in the diagonal matrix
$D=\diag(d_u,d_v)$, and $f,g$ are smooth functions describing the reactions
(production, degradation, interaction). The Neumann condition says that nothing
enters or leaves through the ends of the interval; see \eqref{eq:neumann}.

Ignoring space altogether gives the \emph{kinetic system}
\begin{equation}\label{eq:kinetic-system}
  u' = f(u,v),\qquad v' = g(u,v),
\end{equation}
which describes a well-stirred reactor. Let $(\ub,\vb)$ be a steady state of
\eqref{eq:kinetic-system}, i.e.\ $f(\ub,\vb)=g(\ub,\vb)=0$. The constant functions
$u\equiv\ub$, $v\equiv\vb$ are then also a solution of \eqref{eq:rd-system}: they do
not depend on $t$, their second derivatives vanish, and they satisfy the Neumann
condition. We call $(\ub,\vb)$ a \emph{constant steady state} of the
reaction--diffusion system. It exists for \emph{every} choice of $d_u,d_v$, so we may
ask how its stability depends on diffusion.

The Jacobian matrix of the kinetics at $(\ub,\vb)$ will appear everywhere:
\begin{equation}\label{eq:jacobian-A}
  A = \begin{pmatrix} f_u & f_v\\ g_u & g_v\end{pmatrix},
  \qquad\text{all derivatives evaluated at } (\ub,\vb).
\end{equation}
By \cref{thm:linearized-stability,lem:trace-det}, the steady state is
asymptotically stable for the kinetic system \eqref{eq:kinetic-system} if
\begin{equation*}
  \tr A = f_u+g_v < 0 \qquad\text{and}\qquad \det A = f_ug_v-f_vg_u>0 .
\end{equation*}

\begin{example}\label{ex:turing-schnakenberg-setup}
For the Schnakenberg kinetics \eqref{eq:schnakenberg} with the standard parameters
$a=0.1$, $b=0.9$ we found in \cref{ex:schnakenberg} the steady state
$(\ub,\vb)=(1,0.9)$ and
\begin{equation*}
  A=\begin{pmatrix}0.8 & 1\\ -1.8 & -1\end{pmatrix},\qquad \tr A=-0.2,\qquad \det A=1 .
\end{equation*}
So in a well-stirred reactor this state is stable. We will follow this example
through the whole chapter.
\end{example}

\section{Linearization}\label{sec:turing-linearization}

We perturb the constant steady state, $u=\ub+w_1$, $v=\vb+w_2$, and write
$\bw=(w_1,w_2)^{\top}$. Since $f(\ub,\vb)=0$, Taylor's formula gives
\begin{equation*}
  f(\ub+w_1,\vb+w_2) = f_u w_1 + f_v w_2 + \mathcal{O}(\abs{\bw}^2),
\end{equation*}
and similarly for $g$. Substituting into \eqref{eq:rd-system} and dropping the
quadratic remainder, we obtain the \emph{linearized problem}
\begin{equation}\label{eq:linearized-system}
\begin{aligned}
  \bw_t &= D\,\bw_{xx} + A\,\bw, && x\in(0,L),\ t>0,\\
  \bw_x &= 0, && x\in\{0,L\},\ t>0 .
\end{aligned}
\end{equation}
Written out, $\partial_t w_1 = d_u\,\partial_{xx}w_1 + f_u w_1 + f_v w_2$ and
$\partial_t w_2 = d_v\,\partial_{xx}w_2 + g_u w_1 + g_v w_2$. Note that diffusion
acts on each species separately, while the reactions couple them.

Dropping the remainder is exactly what we do for ODEs, where the principle of
linearized stability (\cref{thm:linearized-stability}) justifies it. For PDEs the
justification is harder; we discuss it in \cref{ch:outlook}. In this chapter we
study the linearized problem only. We measure the size of $\bw$ by
\begin{equation*}
  \norm{\bw(t)}^2 = \norm{w_1(t)}^2+\norm{w_2(t)}^2
  = \int_0^L \abs{\bw(x,t)}^2\dd x .
\end{equation*}

\begin{definition}\label{def:linear-stability-rd}
The constant steady state $(\ub,\vb)$ is \emph{linearly stable} if every solution
of \eqref{eq:linearized-system} (with smooth initial data satisfying the boundary
condition) satisfies $\norm{\bw(t)}\to0$ as $t\to\infty$. It is \emph{linearly
unstable} if \eqref{eq:linearized-system} has a solution with
$\norm{\bw(t)}\to\infty$ as $t\to\infty$.
\end{definition}

If $\bw$ does not depend on $x$, then $\bw_{xx}=0$ and \eqref{eq:linearized-system}
reduces to $\bw'=A\bw$, the linearization of the kinetic system. Perturbations that
are constant in space therefore see only the kinetics. This leads to Turing's
definition.

\begin{definition}[Turing instability]\label{def:ddi}
The constant steady state $(\ub,\vb)$ exhibits \emph{diffusion-driven instability
(DDI)}, also called \emph{Turing instability}, if
\begin{enumerate}[label=(\roman*)]
  \item it is asymptotically stable for the kinetic system, i.e.\ $\tr A<0$ and
        $\det A>0$, and
  \item it is linearly unstable for the reaction--diffusion system
        \eqref{eq:rd-system}.
\end{enumerate}
\end{definition}

In words: the state is stable against spatially homogeneous perturbations, but
some spatially heterogeneous perturbation grows. The instability must be caused by
diffusion, the very process that, on its own, makes everything uniform
(\cref{thm:heat-decay}). That this can happen at all is the surprise of Turing's
paper \cite{Turing1952}.

\section{Normal modes}\label{sec:turing-modes}

For the heat equation, $u(x,t)=e^{-d\lambda_k t}\cos(k\pi x/L)$ solves the problem
because $\cos(k\pi x/L)$ is an eigenfunction of the second derivative with Neumann
boundary conditions: $\varphi_k''=-\lambda_k\varphi_k$, $\varphi_k'(0)=\varphi_k'(L)=0$,
with $\lambda_k=(k\pi/L)^2$ (\cref{prop:neumann-eigen}). For a system we let the
cosine profile carry a \emph{vector} amplitude. It is convenient to introduce the
\emph{mode matrices}
\begin{equation*}
  A_k = A-\lambda_k D
  = \begin{pmatrix} f_u-\lambda_k d_u & f_v\\ g_u & g_v-\lambda_k d_v\end{pmatrix},
  \qquad k=0,1,2,\dots
\end{equation*}
Note that $A_0=A$.

\begin{proposition}[Normal modes]\label{prop:normal-modes}
Let $k\ge0$, let $\mu\in\C$ be an eigenvalue of $A_k$ and let $\bxi\in\C^2\setminus\{0\}$
be a corresponding eigenvector. Then
\begin{equation*}
  \bw(x,t) = e^{\mu t}\cos\Big(\frac{k\pi x}{L}\Big)\,\bxi
\end{equation*}
is a (complex-valued) solution of \eqref{eq:linearized-system}.
\end{proposition}

\begin{proof}
Write $\varphi_k(x)=\cos(k\pi x/L)$. Then $\bw_t=\mu\bw$ and
$\bw_{xx}=e^{\mu t}\varphi_k''\bxi=-\lambda_k\bw$. Hence
\begin{equation*}
  D\bw_{xx}+A\bw = e^{\mu t}\varphi_k\,(A-\lambda_kD)\bxi
  = e^{\mu t}\varphi_k\,\mu\bxi = \bw_t .
\end{equation*}
The boundary condition holds because $\varphi_k'(0)=\varphi_k'(L)=0$.
\end{proof}

So each mode $k$ comes with its own growth rates, namely the eigenvalues of the
$2\times2$ matrix $A_k$. Diffusion enters only through the term $-\lambda_kD$: the
shorter the wavelength (the larger $k$), the stronger the damping. The next theorem
shows that these special solutions are not just examples: every solution is a
superposition of them.

\section{Modal decomposition}\label{sec:turing-decomposition}

We expand a solution in cosines, componentwise, using the coefficients
\eqref{eq:cosine-coefficients}:
\begin{equation*}
  \hatbw_0(t) = \frac1L\int_0^L \bw(x,t)\dd x,\qquad
  \hatbw_k(t) = \frac2L\int_0^L \bw(x,t)\cos\Big(\frac{k\pi x}{L}\Big)\dd x
  \quad (k\ge1).
\end{equation*}
Each $\hatbw_k(t)$ is a vector in $\R^2$.

\begin{theorem}[Modal decomposition]\label{thm:modal-decomposition}
Let $\bw$ be a smooth solution of \eqref{eq:linearized-system}. Then for every
$k\ge0$ the cosine coefficient $\hatbw_k(t)$ solves the linear ODE system
\begin{equation}\label{eq:mode-ode}
  \hatbw_k'(t) = A_k\,\hatbw_k(t),\qquad\text{hence}\qquad
  \hatbw_k(t) = e^{tA_k}\,\hatbw_k(0).
\end{equation}
Consequently
\begin{equation}\label{eq:turing-series}
  \bw(x,t) = \sum_{k=0}^\infty e^{tA_k}\,\hatbw_k(0)\,\cos\Big(\frac{k\pi x}{L}\Big),
\end{equation}
where $\hatbw_k(0)$ are the cosine coefficients of the initial perturbation.
\end{theorem}

\begin{proof}
Let $k\ge1$ and $\varphi_k(x)=\cos(k\pi x/L)$. Since $\bw$ is smooth we may
differentiate under the integral sign and use the equation:
\begin{equation*}
  \hatbw_k'(t) = \frac2L\int_0^L \bw_t\,\varphi_k\dd x
  = D\,\frac2L\int_0^L \bw_{xx}\,\varphi_k\dd x + A\,\frac2L\int_0^L\bw\,\varphi_k\dd x .
\end{equation*}
The last term is $A\hatbw_k$. For the first term we integrate by parts twice,
exactly as in \cref{lem:coefficient-uxx}:
\begin{equation*}
  \int_0^L \bw_{xx}\varphi_k\dd x
  = \big[\bw_x\varphi_k\big]_0^L - \big[\bw\,\varphi_k'\big]_0^L
    + \int_0^L \bw\,\varphi_k''\dd x
  = -\lambda_k\int_0^L\bw\,\varphi_k\dd x .
\end{equation*}
Both boundary terms vanish: the first because $\bw_x=0$ at $x=0,L$, the second
because $\varphi_k'(x)=-\frac{k\pi}{L}\sin(k\pi x/L)$ vanishes at $x=0,L$. Hence
$\hatbw_k' = -\lambda_kD\hatbw_k + A\hatbw_k = A_k\hatbw_k$. For $k=0$ the same
computation with $\varphi_0\equiv1$ gives $\int_0^L\bw_{xx}\dd x=[\bw_x]_0^L=0$, so
$\hatbw_0'=A\hatbw_0=A_0\hatbw_0$. The solution formula for linear ODE systems gives
$\hatbw_k(t)=e^{tA_k}\hatbw_k(0)$, and \eqref{eq:turing-series} is the cosine series
of $\bw(\cdot,t)$ (\cref{thm:cosine-series}).
\end{proof}

\begin{keyidea}
The linearized reaction--diffusion system is equivalent to infinitely many
\emph{decoupled} $2\times2$ linear ODE systems
\begin{equation*}
  \hatbw_k' = (A-\lambda_kD)\,\hatbw_k,\qquad k=0,1,2,\dots,
\end{equation*}
one for each cosine mode. Mode $k$ feels diffusion only through $-\lambda_kD$.
\end{keyidea}

By Parseval's identity \eqref{eq:parseval}, applied to each component,
\begin{equation}\label{eq:turing-parseval}
  \norm{\bw(t)}^2 = L\,\abs{\hatbw_0(t)}^2 + \frac L2\sum_{k=1}^\infty \abs{\hatbw_k(t)}^2 .
\end{equation}
So $\bw$ decays if and only if all the coefficients decay, and the decay must be
uniform in $k$. This is the content of the following criterion.

\begin{theorem}[Spectral criterion]\label{thm:spectral-criterion}
\leavevmode
\begin{enumerate}[label=(\roman*)]
  \item If for some $k\ge0$ the matrix $A_k$ has an eigenvalue $\mu$ with
        $\re\mu>0$, then $(\ub,\vb)$ is linearly unstable.
  \item If for every $k\ge0$ both eigenvalues of $A_k$ have negative real parts,
        then $(\ub,\vb)$ is linearly stable. More precisely, there are constants
        $C\ge1$ and $\gamma>0$ such that every solution of
        \eqref{eq:linearized-system} satisfies
        \begin{equation*}
          \norm{\bw(t)}\le C e^{-\gamma t}\norm{\bw(0)}\qquad\text{for all } t\ge0 .
        \end{equation*}
\end{enumerate}
\end{theorem}

\begin{proof}
(i) Let $\mu=\alpha+i\beta$ with $\alpha>0$, and let $\bxi$ be an eigenvector of
$A_k$. By \cref{prop:normal-modes}, $e^{\mu t}\varphi_k\bxi$ is a solution. Since
the problem is linear with real coefficients, its real part
\begin{equation*}
  \bw(x,t) = e^{\alpha t}\varphi_k(x)\,\big(\cos(\beta t)\,\re\bxi-\sin(\beta t)\,\im\bxi\big)
\end{equation*}
is a real solution. If $\beta=0$ we may take $\bxi$ real, and then
$\norm{\bw(t)}=e^{\alpha t}\norm{\varphi_k}\abs{\bxi}\to\infty$. If $\beta\neq0$,
the vectors $\re\bxi$ and $\im\bxi$ are linearly independent (otherwise $\bxi$
would be a complex multiple of a real vector, and $A_k$, being real, would have a
real eigenvalue with that eigenvector). Hence
$c=\min_{s\in\R}\abs{\cos(s)\re\bxi-\sin(s)\im\bxi}>0$, and
$\norm{\bw(t)}\ge c\,e^{\alpha t}\norm{\varphi_k}\to\infty$.

(ii) \emph{Large $k$.} Let $d_{\min}=\min\{d_u,d_v\}$ and let $\norm{A}$ denote the
operator norm of $A$. For any vector $\mathbf{y}\in\R^2$ we have
$\mathbf{y}\cdot D\mathbf{y}\ge d_{\min}\abs{\mathbf{y}}^2$ and
$\mathbf{y}\cdot A\mathbf{y}\le\norm{A}\abs{\mathbf{y}}^2$. By \eqref{eq:mode-ode},
\begin{equation*}
  \frac{\mathrm d}{\mathrm dt}\abs{\hatbw_k}^2
  = 2\,\hatbw_k\cdot A\hatbw_k - 2\lambda_k\,\hatbw_k\cdot D\hatbw_k
  \le 2\big(\norm{A}-\lambda_kd_{\min}\big)\abs{\hatbw_k}^2 .
\end{equation*}
Since $\lambda_k\to\infty$, there is $K$ such that $\lambda_kd_{\min}\ge\norm{A}+1$ for
all $k\ge K$. For these $k$ we get $\frac{\mathrm d}{\mathrm dt}\abs{\hatbw_k}^2\le
-2\abs{\hatbw_k}^2$, hence $\abs{\hatbw_k(t)}\le e^{-t}\abs{\hatbw_k(0)}$.

\emph{Small $k$.} For each of the finitely many $k<K$ the matrix $A_k$ has
eigenvalues with negative real parts, so by \cref{thm:linear-ode-decay} there are
$C_k\ge1$, $\gamma_k>0$ with $\abs{e^{tA_k}\mathbf{y}}\le C_ke^{-\gamma_kt}\abs{\mathbf{y}}$.

\emph{Conclusion.} Put $C=\max\{C_0,\dots,C_{K-1}\}$ and
$\gamma=\min\{1,\gamma_0,\dots,\gamma_{K-1}\}$. Then
$\abs{\hatbw_k(t)}\le Ce^{-\gamma t}\abs{\hatbw_k(0)}$ for \emph{all} $k\ge0$, and
\eqref{eq:turing-parseval} gives $\norm{\bw(t)}\le Ce^{-\gamma t}\norm{\bw(0)}$.
\end{proof}

The two parts do not cover the borderline case in which some $A_k$ has an
eigenvalue on the imaginary axis and none in the right half-plane. This happens
only for exceptional parameter values (exactly at a bifurcation), and we exclude it.

\begin{remark}[General domains]\label{rem:general-domain}
Nothing in this section used that we work on an interval, except the explicit form
of the eigenpairs $(\lambda_k,\cos(k\pi x/L))$. On a bounded domain
$\Omega\subset\R^n$ one replaces them by the Neumann eigenpairs of $-\Delta$ from
\cref{thm:neumann-eigen-general}, and the same proofs go through. In the language
of operator theory the result says that the spectrum of the linearized operator
$\mathcal L\bw = D\Delta\bw + A\bw$ is
\begin{equation*}
  \sigma(\mathcal L) = \bigcup_{k} \sigma(A-\lambda_kD),
\end{equation*}
the union of the eigenvalues of all mode matrices. This is the form in which the
result usually appears in the literature; see \cite{Perthame2015,Murray2003}.
\end{remark}

\section{When does diffusion destabilize?}\label{sec:turing-criterion}

From now on we assume that the kinetics are stable,
\begin{equation*}
  \tr A<0 \qquad\text{and}\qquad \det A>0,
\end{equation*}
and ask for which $k$ the matrix $A_k$ can have an eigenvalue with positive real
part. Mode $k=0$ is harmless, since $A_0=A$. For $k\ge1$ we use the trace and the
determinant (\cref{lem:trace-det}). The trace is
\begin{equation*}
  \tr A_k = \tr A - \lambda_k(d_u+d_v) < \tr A < 0 .
\end{equation*}
\emph{Diffusion always makes the trace more negative.} The determinant is
\begin{equation*}
  \det A_k = (f_u-\lambda_kd_u)(g_v-\lambda_kd_v) - f_vg_u = h(\lambda_k),
\end{equation*}
where we have introduced the quadratic polynomial
\begin{equation}\label{eq:h-def}
  h(\lambda) = \det(A-\lambda D) = d_ud_v\,\lambda^2 - (d_vf_u+d_ug_v)\,\lambda + \det A .
\end{equation}
By \cref{lem:trace-det}, a real $2\times 2$ matrix with negative trace is stable if
its determinant is positive, and has one positive and one negative real eigenvalue
if its determinant is negative. This gives a complete answer.

\begin{theorem}[Criterion for diffusion-driven instability]\label{thm:ddi-criterion}
Assume $\tr A<0$ and $\det A>0$.
\begin{enumerate}[label=(\alph*)]
  \item If $h(\lambda_k)>0$ for all $k\ge1$, then $(\ub,\vb)$ is linearly stable.
  \item If $h(\lambda_k)<0$ for some $k\ge1$, then $(\ub,\vb)$ is linearly unstable,
        i.e.\ it exhibits diffusion-driven instability. The unstable eigenvalue
        $\mu$ of $A_k$ is real and positive, and the corresponding growing solution
        $e^{\mu t}\cos(k\pi x/L)\,\bxi$ has a real eigenvector $\bxi$.
\end{enumerate}
\end{theorem}

\begin{proof}
(a) $A_0=A$ is stable by assumption, and for $k\ge1$ we have $\tr A_k<0$ and
$\det A_k=h(\lambda_k)>0$, so $A_k$ is stable by \cref{lem:trace-det}. Part (ii) of
\cref{thm:spectral-criterion} applies. (b) If $\det A_k<0$, the eigenvalues of $A_k$
are real with product $\det A_k<0$, so one of them is positive. A real eigenvalue of
a real matrix has a real eigenvector. Part (i) of \cref{thm:spectral-criterion}
applies.
\end{proof}

\begin{keyidea}
Diffusion can destabilize only through the determinant. The trace of every mode
matrix stays negative, so at a Turing instability a \emph{real} eigenvalue becomes
positive: the unstable mode does not oscillate but grows in amplitude while keeping
its spatial shape $\cos(k\pi x/L)$.
\end{keyidea}

\subsection{Conditions that do not depend on the domain}

The polynomial $h$ is an upward parabola with $h(0)=\det A>0$. It takes negative
values for some $\lambda>0$ if and only if its minimum point
$\lambda_{\min}=(d_vf_u+d_ug_v)/(2d_ud_v)$ is positive and the minimum value
$h(\lambda_{\min}) = \det A - (d_vf_u+d_ug_v)^2/(4d_ud_v)$ is negative. (If
$d_vf_u+d_ug_v\le0$, all three terms of $h(\lambda)$ are nonnegative for
$\lambda>0$ and the last one is positive.) Together with the stability of the
kinetics we obtain the four \emph{Turing conditions}:
\begin{equation}\label{eq:turing-conditions}
\begin{aligned}
  &\text{(T1)}\quad f_u+g_v<0, &\qquad
  &\text{(T2)}\quad f_ug_v-f_vg_u>0,\\
  &\text{(T3)}\quad d_vf_u+d_ug_v>0, &\qquad
  &\text{(T4)}\quad (d_vf_u+d_ug_v)^2>4d_ud_v\,(f_ug_v-f_vg_u).
\end{aligned}
\end{equation}
Under (T1)--(T4), $h(\lambda)<0$ exactly on the open interval
$(\lambda_-,\lambda_+)$, where
\begin{equation}\label{eq:band}
  \lambda_\pm = \frac{(d_vf_u+d_ug_v)\pm\sqrt{(d_vf_u+d_ug_v)^2-4d_ud_v\det A}}{2d_ud_v}
  \qquad (0<\lambda_-<\lambda_+).
\end{equation}
We call it the \emph{band of unstable modes}; see \cref{fig:ch05-band}.

\begin{figure}[t]
  \centering
  \includegraphics[width=0.7\textwidth]{ch05-band}
  \caption{The function $h(\lambda)=\det(A-\lambda D)$ for the Schnakenberg kinetics
  with $a=0.1$, $b=0.9$, $d_u=1$ and three values of $\delta=d_v/d_u$. For $\delta<\delta_c$
  the parabola stays positive; at $\delta=\delta_c$ it touches zero; for
  $\delta>\delta_c$ it is negative on the band $(\lambda_-,\lambda_+)$ (shaded for
  $\delta=20$).}
  \label{fig:ch05-band}
\end{figure}

Whether the instability is actually seen depends on the domain, because only the
discrete values $\lambda_k=(k\pi/L)^2$ are available.

\begin{corollary}[Role of the domain size]\label{cor:domain-size}
Assume (T1)--(T4). Then $(\ub,\vb)$ exhibits diffusion-driven instability on
$(0,L)$ if and only if there is an integer $k\ge1$ with
\begin{equation*}
  \sqrt{\lambda_-} < \frac{k\pi}{L} < \sqrt{\lambda_+}
\end{equation*}
(ignoring the borderline case where some $\lambda_k$ equals $\lambda_\pm$). In particular:
\begin{enumerate}[label=(\alph*)]
  \item if $L<\pi/\sqrt{\lambda_+}$, there is no instability: the domain is too small;
  \item if $L\big(\sqrt{\lambda_+}-\sqrt{\lambda_-}\big)>\pi$, there is always an
        instability.
\end{enumerate}
\end{corollary}

\begin{proof}
The first statement is \cref{thm:ddi-criterion} combined with the description of
the set where $h<0$. For (a): $k\ge1$ gives $k\pi/L\ge\pi/L>\sqrt{\lambda_+}$. For
(b): the open interval $\big(L\sqrt{\lambda_-}/\pi,\,L\sqrt{\lambda_+}/\pi\big)$ has
length larger than $1$, so it contains an integer $k$, and $k\ge1$ because the left
endpoint is positive.
\end{proof}

The unstable mode $k$ has wavelength $2L/k=2\pi/q_k$ (\cref{rem:wavenumber}). The
band of unstable wavelengths, $\big(2\pi/\sqrt{\lambda_+},\,2\pi/\sqrt{\lambda_-}\big)$,
does not depend on $L$. Thus the \emph{length scale of the pattern is intrinsic}: it
is set by the kinetics and the diffusion coefficients. The domain only decides which
of the admissible wavelengths fit into it. On a large domain many modes are
unstable (\cref{exer:turing-number-of-modes}).

\subsection{Activators, inhibitors and the diffusion ratio}

The Turing conditions have a clear qualitative meaning.

\begin{corollary}\label{cor:turing-consequences}
Assume (T1)--(T4). Then:
\begin{enumerate}[label=(\alph*)]
  \item $f_u$ and $g_v$ have opposite signs. After renaming the species if
        necessary we may assume $f_u>0>g_v$.
  \item If $f_u>0>g_v$, then $f_vg_u<0$.
  \item If $f_u>0>g_v$, then
        \begin{equation*}
          \delta = \frac{d_v}{d_u} > -\frac{g_v}{f_u} > 1 .
        \end{equation*}
        In particular $d_u\neq d_v$: \emph{with equal diffusion coefficients there is
        no diffusion-driven instability.}
\end{enumerate}
\end{corollary}

\begin{proof}
(a) If $f_u\le0$ and $g_v\le0$, then $d_vf_u+d_ug_v\le0$, contradicting (T3). If
$f_u\ge0$ and $g_v\ge0$, then $f_u+g_v\ge0$, contradicting (T1). (b) By (T2),
$f_vg_u<f_ug_v<0$. (c) Dividing (T3) by $d_uf_u>0$ gives $\delta>-g_v/f_u$, and (T1)
says $f_u<-g_v$, i.e.\ $-g_v/f_u>1$.
\end{proof}

Let us interpret this with $f_u>0>g_v$. The species $u$ enhances its own production
($f_u>0$); it is called the \emph{activator}. The species $v$ limits itself
($g_v<0$). By (c), $v$ must diffuse \emph{faster} than $u$. This is the mechanism of
\emph{local activation and long-range inhibition} \cite{GiererMeinhardt1972}: a small
local excess of $u$ grows by self-enhancement, while the fast-spreading $v$
suppresses growth in the surroundings. Part (b) allows two different interaction
structures:
\begin{itemize}
  \item \emph{activator--inhibitor} kinetics: $f_v<0<g_u$ (the inhibitor $v$ is
        produced by $u$ and suppresses $u$), as in Gierer--Meinhardt
        (\cref{ex:gierer-meinhardt});
  \item \emph{substrate-depletion} kinetics: $f_v>0>g_u$ (the substrate $v$ helps
        to produce $u$ and is consumed by $u$), as in Schnakenberg
        (\cref{ex:schnakenberg}).
\end{itemize}

\begin{remark}[In phase or out of phase]\label{rem:turing-phase}
The two cases produce different patterns. Let $\mu>0$ be the unstable eigenvalue of
$A_k$ with real eigenvector $\bxi=(\xi_1,\xi_2)$. The second row of
$(A_k-\mu)\bxi=0$ reads $g_u\xi_1 + (g_v-\lambda_kd_v-\mu)\xi_2=0$, so
\begin{equation*}
  \xi_2 = \frac{g_u}{\mu+\lambda_kd_v-g_v}\,\xi_1,
  \qquad \mu+\lambda_kd_v-g_v>0 .
\end{equation*}
Hence $\xi_1,\xi_2$ have the sign of $g_u$. For activator--inhibitor kinetics the
peaks of $u$ and $v$ coincide (\emph{in phase}); for substrate depletion the peaks of
$u$ sit at the troughs of $v$ (\emph{out of phase}). Simulations in
\cref{ch:examples} show both cases (\cref{fig:ch06-phase-relation}).
\end{remark}

Conditions (T3) and (T4) involve only the ratio $\delta=d_v/d_u$ (divide them by $d_u$
and $d_u^2$, respectively). They can be combined into a single threshold.

\begin{theorem}[Critical diffusion ratio]\label{thm:critical-ratio}
Assume (T1), (T2) and $f_u>0$. Then (T3) and (T4) hold if and only if
\begin{equation*}
  \delta > \delta_c = \frac{\big(\sqrt{\det A}+\sqrt{-f_vg_u}\,\big)^2}{f_u^2}.
\end{equation*}
At $\delta=\delta_c$ the polynomial $h$ has a double zero at
$\lambda_c=\sqrt{\det A/(d_ud_v)}$.
\end{theorem}

\begin{proof}
By \cref{cor:turing-consequences}(b), $-f_vg_u>0$, so $\delta_c$ is well defined.
Dividing (T4) by $d_u^2$ and expanding, (T4) is equivalent to
\begin{equation*}
  P(\delta) := f_u^2\delta^2 - 2\big(2\det A - f_ug_v\big)\delta + g_v^2 > 0 .
\end{equation*}
The quadratic $P$ has the roots
\begin{equation*}
  \delta_\pm = \frac{(2\det A-f_ug_v)\pm2\sqrt{\det A\,(\det A-f_ug_v)}}{f_u^2}
  = \frac{\big(\sqrt{\det A}\pm\sqrt{-f_vg_u}\,\big)^2}{f_u^2},
\end{equation*}
where we used $\det A-f_ug_v=-f_vg_u$. Thus $\delta_+=\delta_c$, the roots are
distinct, and $P(\delta)>0$ if and only if $\delta<\delta_-$ or $\delta>\delta_c$. By
Vieta's formula $\delta_-\delta_c=g_v^2/f_u^2$, so the geometric mean of the roots is
$-g_v/f_u$ (recall $g_v<0$), and therefore $\delta_-<-g_v/f_u<\delta_c$. Since (T3)
is equivalent to $\delta>-g_v/f_u$, we conclude that (T3) and (T4) together are
equivalent to $\delta>\delta_c$.

At $\delta=\delta_c$ the minimum value of $h$ is zero, so
$h(\lambda)=d_ud_v(\lambda-\lambda_c)^2$ for some $\lambda_c>0$. Comparing the
constant terms gives $d_ud_v\lambda_c^2=\det A$.
\end{proof}

\subsection{Turing's theorem}

Finally, we recover the classical statement: for stable kinetics with an activator,
a sufficiently slowly diffusing activator always produces an instability.

\begin{theorem}\label{thm:turing-small-du}
Assume $\tr A<0$, $\det A>0$ and $f_u>0$. Fix $d_v>0$ and $L>0$. Then there is
$d^*>0$ such that for every $d_u\in(0,d^*)$ the constant steady state $(\ub,\vb)$
exhibits diffusion-driven instability.
\end{theorem}

\begin{proof}
Note first that $g_v<-f_u<0$ by (T1). Since $\lambda_k\to\infty$ and $d_vf_u>0$, we
can choose $k\ge1$ with $\lambda_kd_vf_u>\det A$; put $\eta=\lambda_kd_vf_u-\det A>0$.
Rearranging \eqref{eq:h-def},
\begin{equation*}
  h(\lambda_k) = -\eta + d_u\big(\lambda_k^2d_v - \lambda_kg_v\big),
\end{equation*}
and the bracket is positive and independent of $d_u$. With
$d^*=\eta/(\lambda_k^2d_v-\lambda_kg_v)$ we get $h(\lambda_k)<0$ for all
$d_u\in(0,d^*)$, and \cref{thm:ddi-criterion}(b) applies.
\end{proof}

The theorem says nothing about \emph{nonlinear} instability or about what the
solution does after the perturbation has grown. Both questions are taken up in
\cref{ch:outlook}.

\section{The dispersion relation}\label{sec:turing-dispersion}

It is instructive to forget for a moment that $\lambda$ takes only the discrete
values $\lambda_k$, and to plot the largest real part of the eigenvalues of
$A-\lambda D$ as a function of $\lambda\ge0$:
\begin{equation*}
  \lambda\;\mapsto\;\mu_{\max}(\lambda)=\max\{\re\mu:\ \mu \text{ eigenvalue of } A-\lambda D\}.
\end{equation*}
This curve is called the \emph{dispersion relation}. Each mode $k$ is a point on it,
at $\lambda=\lambda_k$; the mode grows if the point lies above the axis.
\Cref{fig:ch05-dispersion} shows the dispersion relation for our running example.

\begin{figure}[t]
  \centering
  \includegraphics[width=\textwidth]{ch05-dispersion}
  \caption{Dispersion relation for the Schnakenberg kinetics with $a=0.1$, $b=0.9$,
  $d_u=1$. (a) For three values of $\delta=d_v/d_u$. For small $\lambda$ the
  eigenvalues of $A-\lambda D$ are complex, with real part
  $\tfrac12\tr(A-\lambda D)$; the kinks mark the values of $\lambda$ where they
  become real. (b) For $\delta=20$ and $L=40$: the modes $\lambda_k=(k\pi/L)^2$ on
  the curve. Exactly the modes $k=4,\dots,10$ are unstable, and $k=6$ grows fastest.}
  \label{fig:ch05-dispersion}
\end{figure}

As $\delta$ increases through $\delta_c$, the curve first touches the axis at the
single point $\lambda_c$ and then rises above it on the band $(\lambda_-,\lambda_+)$.
In the language of bifurcation theory, $\delta$ is the bifurcation parameter.

\begin{remark}[Turing bifurcation]\label{rem:turing-bifurcation}
At the threshold $\delta=\delta_c$ a single \emph{real} eigenvalue crosses zero, and
it does so at a nonzero wavenumber $q_c=\sqrt{\lambda_c}$. Compare this with a Hopf
bifurcation of an ODE, where a pair of complex eigenvalues crosses the imaginary
axis and oscillations appear. For two species a Hopf-type instability caused by
diffusion is impossible, because the trace of every $A_k$ is negative. With three or
more species it can happen; see \cref{ch:outlook}.
\end{remark}

\begin{example}[Schnakenberg, continued]\label{ex:schnakenberg-turing}
We return to \cref{ex:turing-schnakenberg-setup}: $f_u=0.8$, $f_v=1$, $g_u=-1.8$,
$g_v=-1$, $\tr A=-0.2$, $\det A=1$. Conditions (T1), (T2) hold and $f_u>0$, so
\cref{thm:critical-ratio} gives the critical ratio
\begin{equation*}
  \delta_c = \frac{(1+\sqrt{1.8})^2}{0.64}\approx 8.57 .
\end{equation*}
Take $d_u=1$, $d_v=20$, so $\delta=20>\delta_c$. Then $h(\lambda)=20\lambda^2-15\lambda+1$
and \eqref{eq:band} gives
\begin{equation*}
  \lambda_\pm = \frac{15\pm\sqrt{145}}{40},\qquad
  (\lambda_-,\lambda_+)\approx(0.074,\,0.676),\qquad
  (\sqrt{\lambda_-},\sqrt{\lambda_+})\approx(0.272,\,0.822).
\end{equation*}
By \cref{cor:domain-size}, there is no instability on intervals shorter than
$\pi/\sqrt{\lambda_+}\approx 3.82$. On $(0,40)$ the unstable modes are the integers
$k$ with $0.272<k\pi/40<0.822$, i.e.\ $3.46<k<10.47$: exactly $k=4,\dots,10$. The
largest growth rate, $\mu\approx0.262$, belongs to $k=6$ (\cref{fig:ch05-dispersion}(b)).
So we expect the perturbation to develop, at first, a profile close to
$\cos(6\pi x/40)$, which fits three full wavelengths into the interval. Since $g_u<0$ (substrate depletion), the
peaks of $v$ sit at the troughs of $u$ (\cref{rem:turing-phase}).
\Cref{fig:ch06-1d-evolution} shows a simulation of the full nonlinear system.
\end{example}

\section{Exercises}\label{sec:turing-exercises}

\begin{exercise}\label{exer:turing-equal-diffusion}
Suppose $d_u=d_v=d$. Show directly, without the Turing conditions, that the
eigenvalues of $A_k$ are $\mu-d\lambda_k$, where $\mu$ runs through the eigenvalues
of $A$. Conclude that there is no diffusion-driven instability, and that the same
holds for systems of $n$ species with $D=d\,I$.
\end{exercise}

\begin{exercise}\label{exer:turing-number-of-modes}
Assume (T1)--(T4) and let $N(L)$ be the number of unstable modes on $(0,L)$. Show
that $N(L)$ is finite and that
\begin{equation*}
  \lim_{L\to\infty}\frac{N(L)}{L} = \frac{\sqrt{\lambda_+}-\sqrt{\lambda_-}}{\pi}.
\end{equation*}
\end{exercise}

\begin{exercise}\label{exer:turing-gm}
Consider the Gierer--Meinhardt kinetics \eqref{eq:gierer-meinhardt} with the standard
parameters $a=0$, $b=0.5$ (\cref{ex:gierer-meinhardt}).
\begin{enumerate}[label=(\alph*)]
  \item Verify (T1), (T2) and show that $\delta_c=(\sqrt{0.5}+1)^2/0.25\approx 11.66$.
  \item For $d_u=1$, $d_v=30$ compute the band $(\lambda_-,\lambda_+)$, the shortest
        interval length that allows an instability, and the unstable modes on $(0,40)$.
  \item Is the pattern in phase or out of phase?
\end{enumerate}
\end{exercise}

\begin{exercise}\label{exer:turing-lengyel-epstein}
For the Lengyel--Epstein kinetics \eqref{eq:lengyel-epstein} with $a=10$ show that
$(\ub,\vb)=(2,5)$, that the kinetics are stable if and only if $b>3.5$, and that
\begin{equation*}
  \delta_c = \frac{25}{49}\big(\sqrt2+1.6\big)^2\,b\approx 4.64\,b .
\end{equation*}
\end{exercise}

\begin{exercise}[Only two numbers matter]\label{exer:turing-scaling}
Let $\ell>0$ and set $\tilde x=x/\ell$. Show that $(u,v)$ solves \eqref{eq:rd-system}
on $(0,L)$ if and only if $\tilde u(\tilde x,t)=u(\ell\tilde x,t)$,
$\tilde v(\tilde x,t)=v(\ell\tilde x,t)$ solve the same system on $(0,L/\ell)$ with
diffusion coefficients $d_u/\ell^2$ and $d_v/\ell^2$. Deduce that, for given kinetics,
the linear stability of $(\ub,\vb)$ depends only on $\delta=d_v/d_u$ and on $L^2/d_u$.
Reformulate \cref{thm:turing-small-du} as a statement about large domains.
\end{exercise}

\begin{exercise}[$\ast$]\label{exer:turing-explicit-roots}
Instead of the trace--determinant argument one can compute the eigenvalues of
$A_k$ explicitly:
\begin{equation*}
  \mu_{1,2} = -p_k\pm\sqrt{p_k^2-h(\lambda_k)},\qquad
  p_k=\tfrac12\big(\lambda_k(d_u+d_v)-\tr A\big).
\end{equation*}
Verify this formula, and use it to give a second proof of \cref{thm:ddi-criterion}.
Show in addition that $\re\mu_2<0$ for every $k$, and that $\max\{\re\mu: \mu\in\sigma(A_k),\ k\ge0\}$
is attained. Why is the last statement useful?
\end{exercise}
